% TOP_PROBLEM: {"id": 3096, "title": "The Hodge Conjecture", "category_id": 6, "category": {"id": 6, "name": "geometry", "display_name": "Geometry", "description": "Euclidean and non-Euclidean geometry, geometric structures.", "slug": "geometry", "order_index": 6, "created_at": "2026-07-31T15:26:25.671Z"}, "status": "open", "problem_number": "OPG-1803", "set_id": 12, "statusQualification": "The standard rational Hodge conjecture is used, explicitly restoring smoothness of the ambient projective variety."}
% FIRST_POSTED: 2026-10-06
% ATTEMPT_STATUS: unresolved
% OPENAI_CLAIM: {"schema_version":1,"source_repo":"https://github.com/openai/math","source_commit":"adc7f1241b42e322a6451854ab7e4b4c146bf78a","release_date":"2026-10-06","match":"partial","resolution":"partial","independently_reviewed":false,"second_pass":"n/a","families":[{"family":"032","title":"Hodge and Kuga–Satake results for all projective K3 surfaces","subject":"Algebraic and complex geometry","manuscripts":[{"title":"A Conditional Reduction for Algebraic Kuga–Satake Correspondences","pdf_path":"preprints/A-Conditional-Reduction-for-Algebraic-Kuga-Satake-Correspondences-September-10-2026/paper.pdf","readme_path":"preprints/A-Conditional-Reduction-for-Algebraic-Kuga-Satake-Correspondences-September-10-2026/README.md","date":"2026-09-10","theorem_ref":"main:generic (sections/introduction.tex:36)","statement_sources":[{"path":"preprints/A-Conditional-Reduction-for-Algebraic-Kuga-Satake-Correspondences-September-10-2026/build/sections/introduction.tex","line":36,"label":"main:generic"}]},{"title":"Abelian covers, Gale correspondences, and the Hodge conjecture for powers","pdf_path":"preprints/Abelian-covers-Gale-correspondences-and-the-Hodge-conjecture-for-powers-September-30-2026/paper.pdf","readme_path":"preprints/Abelian-covers-Gale-correspondences-and-the-Hodge-conjecture-for-powers-September-30-2026/README.md","date":"2026-09-30","theorem_ref":"thm:main (paper.tex:146); p4:abelian-covers (paper.tex:104)","statement_sources":[{"path":"preprints/Abelian-covers-Gale-correspondences-and-the-Hodge-conjecture-for-powers-September-30-2026/build/paper.tex","line":146,"label":"thm:main"},{"path":"preprints/Abelian-covers-Gale-correspondences-and-the-Hodge-conjecture-for-powers-September-30-2026/build/paper.tex","line":104,"label":"p4:abelian-covers"}]},{"title":"Algebraic Kuga–Satake correspondences and Hodge conjectures on a K3 quadratic locus","pdf_path":"preprints/Algebraic-Kuga-Satake-correspondences-and-Hodge-conjectures-on-a-K3-quadratic-locus-September-30-2026/paper.pdf","readme_path":"preprints/Algebraic-Kuga-Satake-correspondences-and-Hodge-conjectures-on-a-K3-quadratic-locus-September-30-2026/README.md","date":"2026-09-30","theorem_ref":"main:theorem (paper.tex:179)","statement_sources":[{"path":"preprints/Algebraic-Kuga-Satake-correspondences-and-Hodge-conjectures-on-a-K3-quadratic-locus-September-30-2026/build/paper.tex","line":179,"label":"main:theorem"}]},{"title":"Algebraicity of Kuga–Satake Correspondences for K3 Surfaces","pdf_path":"preprints/Algebraicity-of-Kuga-Satake-Correspondences-for-K3-Surfaces-October-3-2026/manuscript.pdf","readme_path":"preprints/Algebraicity-of-Kuga-Satake-Correspondences-for-K3-Surfaces-October-3-2026/README.md","date":"2026-10-03","theorem_ref":"thm:main (manuscript.tex:109)","statement_sources":[{"path":"preprints/Algebraicity-of-Kuga-Satake-Correspondences-for-K3-Surfaces-October-3-2026/build/manuscript.tex","line":109,"label":"thm:main"}]},{"title":"Algebraicity of Weil classes on split abelian eightfolds","pdf_path":"preprints/Algebraicity-of-Weil-classes-on-split-abelian-eightfolds-September-18-2026/paper.pdf","readme_path":"preprints/Algebraicity-of-Weil-classes-on-split-abelian-eightfolds-September-18-2026/README.md","date":"2026-09-18","theorem_ref":"intro:main (sections/01-introduction.tex:51)","statement_sources":[{"path":"preprints/Algebraicity-of-Weil-classes-on-split-abelian-eightfolds-September-18-2026/build/sections/01-introduction.tex","line":51,"label":"intro:main"}]},{"title":"The rational Hodge conjecture for CM abelian varieties","pdf_path":"preprints/The-rational-Hodge-conjecture-for-CM-abelian-varieties-September-30-2026/paper.pdf","readme_path":"preprints/The-rational-Hodge-conjecture-for-CM-abelian-varieties-September-30-2026/README.md","date":"2026-09-30","theorem_ref":"thm:cm (sections/01-introduction.tex:17); cc:ghc (sections/07-consequences.tex:154); cc:tate (sections/07-consequences.tex:213); cc:standard (sections/07-consequences.tex:274)","statement_sources":[{"path":"preprints/The-rational-Hodge-conjecture-for-CM-abelian-varieties-September-30-2026/build/sections/01-introduction.tex","line":17,"label":"thm:cm"},{"path":"preprints/The-rational-Hodge-conjecture-for-CM-abelian-varieties-September-30-2026/build/sections/07-consequences.tex","line":154,"label":"cc:ghc"},{"path":"preprints/The-rational-Hodge-conjecture-for-CM-abelian-varieties-September-30-2026/build/sections/07-consequences.tex","line":213,"label":"cc:tate"},{"path":"preprints/The-rational-Hodge-conjecture-for-CM-abelian-varieties-September-30-2026/build/sections/07-consequences.tex","line":274,"label":"cc:standard"}]},{"title":"The rational Hodge conjecture for products of K3 surfaces","pdf_path":"preprints/The-rational-Hodge-conjecture-for-products-of-K3-surfaces-October-4-2026/hodge-conjecture-products-k3.pdf","readme_path":"preprints/The-rational-Hodge-conjecture-for-products-of-K3-surfaces-October-4-2026/README.md","date":"2026-10-04","theorem_ref":"thm:main (manuscript/introduction.tex:14)","statement_sources":[{"path":"preprints/The-rational-Hodge-conjecture-for-products-of-K3-surfaces-October-4-2026/build/manuscript/introduction.tex","line":14,"label":"thm:main"}]},{"title":"Weil classes and Hodge classes on abelian powers","pdf_path":"preprints/Weil-classes-and-Hodge-classes-on-abelian-powers-September-30-2026/paper.pdf","readme_path":"preprints/Weil-classes-and-Hodge-classes-on-abelian-powers-September-30-2026/README.md","date":"2026-09-30","theorem_ref":"p3ab:split-sixfolds (paper.tex:79); p3ab:quadratic-five (paper.tex:88)","statement_sources":[{"path":"preprints/Weil-classes-and-Hodge-classes-on-abelian-powers-September-30-2026/build/paper.tex","line":79,"label":"p3ab:split-sixfolds"},{"path":"preprints/Weil-classes-and-Hodge-classes-on-abelian-powers-September-30-2026/build/paper.tex","line":88,"label":"p3ab:quadratic-five"}]}],"lean":null,"reasoning_trace":null},{"family":"040","title":"Bloch’s conjecture for complex surfaces","subject":"Algebraic and complex geometry","manuscripts":[{"title":"Bloch's conjecture for surfaces with p_g=0","pdf_path":"preprints/Blochs-Conjecture-for-Surfaces-with-pg-equals-q-equals-0-September-24-2026/paper.pdf","readme_path":"preprints/Blochs-Conjecture-for-Surfaces-with-pg-equals-q-equals-0-September-24-2026/README.md","date":"2026-09-24","theorem_ref":"cor:tate-motives (sections/introduction.tex:234); thm:bloch (sections/introduction.tex:18); cor:bloch-pg-zero (sections/introduction.tex:31)","statement_sources":[{"path":"preprints/Blochs-Conjecture-for-Surfaces-with-pg-equals-q-equals-0-September-24-2026/build/sections/introduction.tex","line":234,"label":"cor:tate-motives"},{"path":"preprints/Blochs-Conjecture-for-Surfaces-with-pg-equals-q-equals-0-September-24-2026/build/sections/introduction.tex","line":18,"label":"thm:bloch"},{"path":"preprints/Blochs-Conjecture-for-Surfaces-with-pg-equals-q-equals-0-September-24-2026/build/sections/introduction.tex","line":31,"label":"cor:bloch-pg-zero"}]}],"lean":null,"reasoning_trace":null}]}
\documentclass[11pt]{article}
\usepackage[margin=25mm]{geometry}
\usepackage{fontspec}
\setmainfont{Latin Modern Roman}
\usepackage{amsmath,amssymb,mathtools}
\usepackage{unicode-math}
\setmathfont{Latin Modern Math}
\usepackage{enumitem}
\usepackage{xurl,hyperref,bookmark}
\hypersetup{colorlinks=true,urlcolor=blue}
\setcounter{secnumdepth}{0}
\setlength{\parindent}{0pt}
\setlength{\parskip}{5pt}
\setlength{\emergencystretch}{3em}
\begin{document}
\section*{The Hodge Conjecture}

\subsection{OpenAI's claim}
Manuscript 1 (family 032, 2026-09-10): A nonzero algebraic map from generic primitive K3 cohomology to an abelian target is claimed to recover the exact Kuga--Satake correspondence; existence of that starting map remains an assumption; Manuscript 2 (family 032, 2026-09-30): For tensor-Hodge-generic abelian-cover Jacobians and the specified very general diagonal complete intersections, OpenAI claims the rational Hodge assertion on every self-power; Manuscript 3 (family 032, 2026-09-30): On the specified quadratic-lattice locus of projective K3 surfaces, all rational Hodge classes on every power are claimed to be algebraic; an exact algebraic Kuga--Satake tensor also gives a conditional criterion beyond that locus. Manuscript 4 (family 032, 2026-10-03): For every projective complex K3 surface, the precisely normalized Kuga--Satake embedding into the full even-Clifford abelian target is claimed to come from an algebraic cycle; Manuscript 5 (family 032, 2026-09-18): Every rational Weil class of a split eight-dimensional abelian variety is claimed to have an algebraic representative, uniformly over imaginary-quadratic fields and compatible polarizations; Manuscript 6 (family 032, 2026-09-30): The manuscripts claim algebraicity of rational Hodge classes on CM abelian varieties, with consequences for generalized Hodge classes and finite-field Tate classes. Manuscript 7 (family 032, 2026-10-04): Rational Hodge classes on any finite product of projective complex K3 surfaces are claimed to be algebraic, even when the factors have unrelated periods or extra endomorphisms; Manuscript 8 (family 032, 2026-09-30): Algebraicity is claimed for Hodge classes on powers of split Weil sixfolds and of abelian varieties of dimension at most five carrying an imaginary-quadratic action, including special periods; Manuscript 9 (family 040, 2026-09-24): For every smooth complex projective surface with $p_g=0$, OpenAI claims that the Albanese map on integral degree-zero cycles is an isomorphism.

\subsection{Mathematical statement of the claimed result}
\paragraph{Mathematical claim 1.} Let $(S,h)$ be a primitively amply polarized projective complex K3 surface, put $V=h^\perp\subset H^2(S,\mathbb Q)$ with cup-product form $q$, and assume $\operatorname{Hg}(V)=\operatorname{SO}(V,q)$. This hypothesis implies $T(S)=\operatorname{NS}(S)_{\mathbb Q}^{\perp}=V$. For any stipulated rational quadratic polarization and full standard even-Clifford Kuga--Satake construction, including parity, weight-one polarization and isogenous realization, write $A_S$ for its abelian variety and $\kappa_S:V\to H^2(A_S^2,\mathbb Q)$ for the prescribed map. OpenAI asserts equivalence of these three existence statements: (i) some abelian variety $B$ admits a rational codimension-two cycle on $S\times B$ inducing a nonzero map from $V$ to $H^2(B,\mathbb Q)$; (ii) a smooth connected projective surface $Y$ admits a dominant generically finite map $f:Y\to S$ with
\[
f_*\operatorname{im}\bigl(\bigwedge^2 H^0(Y,\Omega_Y^1)\longrightarrow H^0(Y,\Omega_Y^2)\bigr)\ne0;
\]
(iii) some $\Gamma\in\operatorname{CH}^2(S\times A_S^2)_{\mathbb Q}$ satisfies $\Gamma_*|_V=\kappa_S$. Thus algebraicity is conditional on either equivalent geometric input; this theorem does not assert that every K3 surface has that input.

\paragraph{Mathematical claim 2.} Fix $n\ge2$, $d\ge0$, $c\in\{0,1,2\}$ and $m=d+c+1$. For a rank-$c$ complex matrix $A=(a_{\alpha i})$ with every maximal coordinate minor nonzero, define the smooth diagonal complete intersection
\[
X_A=\{[x_1:\cdots:x_m]\in\mathbb P^{m-1}:\ \sum_i a_{\alpha i}x_i^n=0\quad(1\le\alpha\le c)\}.
\]
OpenAI asserts surjectivity of $\operatorname{CH}^p(X_A^u)_{\mathbb Q}\to H^{2p}(X_A^u,\mathbb Q)\cap H^{p,p}(X_A^u)$ for every $u\ge1$ and every $p$, outside a union, indexed by a countable set, of proper algebraic closed subsets of this smooth parameter space. The exceptional set is unnecessary when $c\le1$ or $d=0$.
A second assertion concerns a fixed connected finite abelian topological cover of a compact oriented surface with specified branch monodromy, allowing unbranched and trivial covers. Vary the full marked complex structure and branch points. At any parameter $s$ for which each rational tensor of type $(a,a)$ in $H^1(C_s,\mathbb Q)^{\otimes 2a}$ keeps its Hodge type at every point of the marked family, the same cycle-class surjectivity holds on
\[
(\operatorname{Jac}(C_s)\times M)^N\quad(N\ge1),\qquad C_s^a\times M^b\quad(a,b\ge0),
\]
for every complex CM abelian variety $M$. CM means that its degree-one Hodge group is a torus. This is a Hodge-generic assertion, with the generic locus chosen independently of $M$.

\paragraph{Mathematical claim 3.} For a projective complex K3 surface $S$, write $T_S=\operatorname{NS}(S)_{\mathbb Q}^{\perp}$ with cup-product quadratic form $q$. Assume a rational isometric embedding
\[
(T_S,q)\hookrightarrow\mathbb U_{\mathbb Q}^{\oplus2}\perp\langle-1\rangle^4,
\]
where $\mathbb U$ is the hyperbolic plane with Gram matrix $\left(\begin{smallmatrix}0&1\\1&0\end{smallmatrix}\right)$. Choose any full even-Clifford realization $H^1(A,\mathbb Q)=C^+(T_S,q)$, any $w\in T_S$ with $q(w)\ne0$, and any rational weight-one polarization. The Kuga--Satake map is formed from $v\mapsto(a\mapsto vaw)$ and this polarization. OpenAI asserts existence of $\Gamma\in\operatorname{CH}^2(S\times A^2)_{\mathbb Q}$ inducing this exact prescribed map on $T_S$, also after transporting it to an isogenous realization. It further claims rational Hodge and generalized Hodge conjectures in all degrees of $S^m$ for every $m\ge1$. The quadratic embedding is a hypothesis, so this theorem alone is restricted to that locus.

\paragraph{Mathematical claim 4.} Let $S$ be any smooth projective complex K3 surface and let $T(S)=\operatorname{NS}(S)_{\mathbb Q}^{\perp}$ in $H^2(S,\mathbb Q)$. Fix the full even-Clifford Kuga--Satake construction, including the rational quadratic polarization, parity identification and rational weight-one polarization, with associated abelian variety $A_S$. These choices specify $\kappa_S:T(S)\to H^2(A_S^2,\mathbb Q)$. OpenAI claims
\[
\exists\Gamma_S\in\operatorname{CH}^2(S\times A_S^2)_{\mathbb Q}\quad
(p_{A_S^2})_*\bigl(\operatorname{cl}(\Gamma_S)\smile p_S^*x\bigr)=\kappa_S(x)
\quad\text{for all }x\in T(S).
\]
Here $\operatorname{CH}^2$ is the group of rational codimension-two cycles modulo rational equivalence and $p_S,p_{A_S^2}$ are the projections. The assertion includes every isogenous realization with the prescribed map transported by its rational Hodge isomorphism.

\paragraph{Mathematical claim 5.} Let $K=\mathbb Q(\sqrt{-d})$ be imaginary quadratic, $A$ a complex abelian variety of dimension eight with a $K$-action, and $\lambda$ a compatible polarization of Weil type, with eigenspace multiplicities $(4,4)$. Put $V=H_1(A,\mathbb Q)$, let $\psi$ be the rational Riemann form and $J$ multiplication by $\sqrt{-d}$. Assume the associated Hermitian form
\[
h(x,y)=\psi(Jx,y)+\sqrt{-d}\,\psi(x,y)
\]
vanishes on a four-dimensional $K$-subspace of $V$ (the split condition). Define the rational Weil space $W_K(A)=\bigwedge_K^8H^1(A,\mathbb Q)\subset H^8(A,\mathbb Q)$ using the two conjugate $K$-eigenspace exterior powers. OpenAI claims
\[
W_K(A)\subset\operatorname{im}\bigl(\operatorname{CH}^4(A)_{\mathbb Q}\xrightarrow{\operatorname{cl}}H^8(A,\mathbb Q)\bigr).
\]
The field and compatible polarization type are arbitrary; extra endomorphisms are allowed. The statement concerns these Weil classes, rather than all Hodge classes on all varieties.

\paragraph{Mathematical claim 6.} For a smooth connected projective complex surface $S$ put $p_g=h^0(S,\omega_S)$ and $q=h^1(S,\mathcal O_S)$. OpenAI asserts the integral statements
\[
p_g=q=0\ \Longrightarrow\ \deg:\operatorname{CH}_0(S)\xrightarrow{\sim}\mathbb Z,
\qquad
p_g=0\ \Longrightarrow\ \operatorname{alb}_S:\operatorname{CH}_0(S)^0\xrightarrow{\sim}\operatorname{Alb}(S)(\mathbb C),
\]
where $\operatorname{CH}_0(S)^0$ is the subgroup of degree-zero zero-cycles modulo rational equivalence. Under $p_g=0$, it also claims that the rational Chow motive is finite-dimensional in Kimura's sense and has zero transcendental degree-two summand. When $p_g=q=0$ it claims
\[
h(S)_{\mathbb Q}\simeq\mathbf1\oplus\mathbb L^{\oplus b_2(S)}\oplus\mathbb L^2,
\]
where $\mathbb L$ is the Lefschetz motive. Under the joint hypothesis $p_g=q=0$, for each power $Y=S^m$ and each Hilbert scheme $Y=S^{[n]}$, the corresponding rational motive is pure Tate and $\operatorname{CH}^i(Y)_{\mathbb Q}\xrightarrow{\sim}H^{2i}(Y,\mathbb Q)$, with $H^{2i+1}(Y,\mathbb Q)=0$, for every $i\ge0$.

\paragraph{Mathematical claim 7.} A complex abelian variety $A$ is CM if $\operatorname{End}(A)\otimes\mathbb Q$ contains a commutative semisimple algebra of dimension $2\dim A$. For every such $A$ and every $p\ge0$, OpenAI claims
\[
\operatorname{CH}^p(A)_{\mathbb Q}\longrightarrow H^{2p}(A,\mathbb Q)\cap H^{p,p}(A)
\quad\text{is surjective}.
\]
This also applies to arbitrary finite products of CM abelian varieties, including their powers. If $W\subset H^k(A,\mathbb Q)$ is a rational Hodge substructure and $W(r)$ is effective (all Hodge bidegrees are nonnegative), it claims a closed algebraic $Z\subset A$ of codimension at least $r$ with $W\subset\ker[H^k(A,\mathbb Q)\to H^k(A\setminus Z,\mathbb Q)]$.
The manuscript additionally asserts that, for every abelian variety $A/\mathbb F_q$, every $r\ge0$ and prime $\ell\ne\operatorname{char}\mathbb F_q$,
\[
\operatorname{CH}^r(A)\otimes\mathbb Q_\ell\longrightarrow
H^{2r}_{\mathrm{\acute et}}(A_{\overline{\mathbb F}_q},\mathbb Q_\ell(r))^{\operatorname{Gal}(\overline{\mathbb F}_q/\mathbb F_q)}
\quad\text{is surjective}.
\]
Numerical and $\ell$-adic homological equivalence are claimed to agree on rational cycles. Over any algebraically closed field, for an abelian $g$-fold and rational ample class $h$, the form $(-1)^r\deg(xyh^{g-2r})$ is claimed positive definite on the real extension of $\ker[h^{g-2r+1}:A^r_{\mathrm{num}}\to A^{g-r+1}_{\mathrm{num}}]$ for $0\le r\le g/2$.

\paragraph{Mathematical claim 8.} Choose any integer $m\ge1$ and any projective complex K3 surfaces $S_1,\ldots,S_m$, allowing repeated factors. For $X=\prod_{j=1}^m S_j$, OpenAI claims
\[
\operatorname{CH}^p(X)_{\mathbb Q}\xrightarrow{\operatorname{cl}}
H^{2p}(X,\mathbb Q)\cap H^{p,p}(X)
\quad\text{is surjective for every }0\le p\le2m.
\]
There are no Picard-number, period or endomorphism-field restrictions. The quantified varieties are products of K3 surfaces, so this statement has smaller scope than the Hodge conjecture for arbitrary smooth projective complex varieties.

\paragraph{Mathematical claim 9.} Let $E=\mathbb Q(u)$ be imaginary quadratic with $u^2=-d$, $d>0$, and let $A$ be a complex abelian variety with a unital $E$-action. OpenAI claims rational cycle-class surjectivity
\[
\operatorname{CH}^p(A^m)_{\mathbb Q}\longrightarrow H^{2p}(A^m,\mathbb Q)\cap H^{p,p}(A^m)
\quad(m\ge1,\ p\ge0)
\]
in either of these two cases: $\dim A\le5$; or $\dim A=6$ and the rational Hermitian form associated to a compatible polarization is hyperbolic with signature $(3,3)$. In the second case, with rational Riemann form $e$, this means $h_E(x,y)=e(x,uy)+u e(x,y)$ on $H_1(A,\mathbb Q)$ has a totally isotropic three-dimensional $E$-subspace. Compatibility means $e(ax,y)=e(x,\bar a y)$ for all $a\in E$.

\subsection{Outline of the argument}
Manuscript 1 (family 032, 2026-09-10): Acting independently on each weight-one factor reconstructs the Clifford multiplicities, and rational descent treats nonsplit algebras. Hilbert parameter spaces carry the specified flat cycle class to special periods, while signed square-root comparisons keep the polarization normalization fixed. Manuscript 2 (family 032, 2026-09-30): Deck-character projections and equivariant period calculations identify the necessary monodromy invariants of curve covers. Correspondences dominated by CM sources supply determinant tensors, while Gale duality transfers the character-level construction to diagonal complete intersections with the required supports and degree shifts. Manuscript 3 (family 032, 2026-09-30): Clifford multiplication carries exterior powers of the transcendental space into degree-two cohomology of an abelian target. Transposed cycles and a polynomial inverse on the finite-dimensional image recover these summands, so divisor algebraicity supplies the required classes without an inverse Lefschetz correspondence.

Manuscript 4 (family 032, 2026-10-03): The construction engineers a Lagrangian whose second cohomology has a nineteen-dimensional deficit, yielding a corresponding bound on self-extensions after mirror realization. The resulting characteristic class is transported through K3 moduli by semiregularity and specialization, retaining the prescribed Clifford tensor. Manuscript 5 (family 032, 2026-09-18): An embedded spin Lagrangian detects the exceptional middle-cohomology component and makes two homological obstruction maps injective. Finite-cutoff weighted Floer operations are assembled into algebraic module identities; cancellation in alternating specialization then transports the desired Chern character to an algebraic cycle. Manuscript 6 (family 032, 2026-09-30): The argument realizes four-factor balanced tensors using unitary theta constructions. Elementary switches between equal-sum lists, followed by algebraic polarization contractions and descent of coefficients, build arbitrary balanced CM tensors from these elementary cycle classes.

Manuscript 7 (family 032, 2026-10-04): A Lagrangian immersion is used to produce prescribed Euler pairings while bounding the second self-extension space of the associated complex. Clifford-algebra calculations then identify its transcendental cohomology component; algebraic Kuga--Satake tensors and their contractions supply the classes needed on products. Manuscript 8 (family 032, 2026-09-30): A coherent seed and its Fourier--Mukai transform produce the split-sixfold Weil plane through an obstruction calculation and deformation. A determinant description of the actual Hodge group, combined with the CM result, extends these algebraic generators to powers and to the stated low-dimensional cases. Manuscript 9 (family 040, 2026-09-24): A determinant-constrained Quot obstruction theory and torus localization produce a Chow relation involving the diagonal. A nested-Hilbert recursion leaves two surface factors unintegrated, retaining an actual correspondence identity; its external-product terms annihilate degree-zero cycles, and the diagonal coefficient gives the required vanishing after the standard reductions.

\subsection{Scope relative to this problem}
Manuscript 1 (family 032, 2026-09-10): A conditional geometric criterion for algebraicity of the Kuga–Satake tensor; this problem quantify over every smooth projective complex variety, so the restricted object class or conditional tensor criterion leaves the general Hodge assertion open.

Manuscript 2 (family 032, 2026-09-30): Hodge classes on powers of specified Hodge-generic abelian-cover Jacobians and diagonal complete intersections; this problem quantify over every smooth projective complex variety, so the restricted object class or conditional tensor criterion leaves the general Hodge assertion open.

Manuscript 3 (family 032, 2026-09-30): Hodge and generalized Hodge for powers of K3 surfaces whose transcendental lattice has the specified quadratic embedding; this problem quantify over every smooth projective complex variety, so the restricted object class or conditional tensor criterion leaves the general Hodge assertion open.

Manuscript 4 (family 032, 2026-10-03): Algebraicity of the prescribed Kuga–Satake tensor for every projective complex K3 surface; this problem quantify over every smooth projective complex variety, so the restricted object class or conditional tensor criterion leaves the general Hodge assertion open.

Manuscript 5 (family 032, 2026-09-18): Algebraicity of Weil classes for split Weil-type abelian eightfolds; this problem quantify over every smooth projective complex variety, so the restricted object class or conditional tensor criterion leaves the general Hodge assertion open.

Manuscript 6 (family 032, 2026-09-30): Rational Hodge for complex CM abelian varieties, with the stated CM generalized-Hodge and finite-field abelian consequences; this problem quantify over every smooth projective complex variety, so the restricted object class or conditional tensor criterion leaves the general Hodge assertion open.

Manuscript 7 (family 032, 2026-10-04): Rational cycle-class surjectivity on every product of projective complex K3 surfaces; this problem quantify over every smooth projective complex variety, so the restricted object class or conditional tensor criterion leaves the general Hodge assertion open.

Manuscript 8 (family 032, 2026-09-30): Hodge classes on powers of split Weil sixfolds and quadratic-action abelian varieties of dimension at most five; this problem quantify over every smooth projective complex variety, so the restricted object class or conditional tensor criterion leaves the general Hodge assertion open.

Manuscript 9 (family 040, 2026-09-24): The pure-Tate motive corollary proves Hodge for powers and Hilbert schemes of complex surfaces with p\_g=q=0. It leaves the universal smooth projective Hodge statement untouched.

\subsection{Formal verification}
Family 032 has no Lean scope page in this snapshot; no formalization of these manuscript results is listed.

Family 040 has no Lean scope page in this snapshot; no formalization of these manuscript results is listed.

This site has not executed the Lean files or checked their dependencies. These descriptions report the scope of the linked declarations, rather than a new verification of the proofs.

\subsection{Status caveat}
OpenAI attributes these results to an internal model. Its release notice cautions that results lacking full formalization can contain mistakes. This site has not independently assessed the proofs; a solved label records the scope of the reported claim.

\subsection{Sources}
\begin{itemize}
\item OpenAI, \emph{A Conditional Reduction for Algebraic Kuga–Satake Correspondences}, 2026-09-10; citation key \texttt{OAI:A-Conditional-Reduction-for-Algebraic-Kuga-Satake-Correspondences-September-10-2026}. \href{https://github.com/openai/math/blob/adc7f1241b42e322a6451854ab7e4b4c146bf78a/preprints/A-Conditional-Reduction-for-Algebraic-Kuga-Satake-Correspondences-September-10-2026/paper.pdf}{PDF (GitHub)}; \href{https://github.com/openai/math/raw/adc7f1241b42e322a6451854ab7e4b4c146bf78a/preprints/A-Conditional-Reduction-for-Algebraic-Kuga-Satake-Correspondences-September-10-2026/paper.pdf}{PDF (direct)}; \href{https://github.com/openai/math/blob/adc7f1241b42e322a6451854ab7e4b4c146bf78a/preprints/A-Conditional-Reduction-for-Algebraic-Kuga-Satake-Correspondences-September-10-2026/README.md}{Manuscript page and citation}; \href{https://github.com/openai/math/blob/main/preprints/A-Conditional-Reduction-for-Algebraic-Kuga-Satake-Correspondences-September-10-2026/paper.pdf}{Latest version}; Statement source: \href{https://github.com/openai/math/blob/adc7f1241b42e322a6451854ab7e4b4c146bf78a/preprints/A-Conditional-Reduction-for-Algebraic-Kuga-Satake-Correspondences-September-10-2026/build/sections/introduction.tex\#L36}{main:generic}.
\item OpenAI, \emph{Abelian covers, Gale correspondences, and the Hodge conjecture for powers}, 2026-09-30; citation key \texttt{OAI:Abelian-covers-Gale-correspondences-and-the-Hodge-conjecture-for-powers-September-30-2026}. \href{https://github.com/openai/math/blob/adc7f1241b42e322a6451854ab7e4b4c146bf78a/preprints/Abelian-covers-Gale-correspondences-and-the-Hodge-conjecture-for-powers-September-30-2026/paper.pdf}{PDF (GitHub)}; \href{https://github.com/openai/math/raw/adc7f1241b42e322a6451854ab7e4b4c146bf78a/preprints/Abelian-covers-Gale-correspondences-and-the-Hodge-conjecture-for-powers-September-30-2026/paper.pdf}{PDF (direct)}; \href{https://github.com/openai/math/blob/adc7f1241b42e322a6451854ab7e4b4c146bf78a/preprints/Abelian-covers-Gale-correspondences-and-the-Hodge-conjecture-for-powers-September-30-2026/README.md}{Manuscript page and citation}; \href{https://github.com/openai/math/blob/main/preprints/Abelian-covers-Gale-correspondences-and-the-Hodge-conjecture-for-powers-September-30-2026/paper.pdf}{Latest version}; Statement source: \href{https://github.com/openai/math/blob/adc7f1241b42e322a6451854ab7e4b4c146bf78a/preprints/Abelian-covers-Gale-correspondences-and-the-Hodge-conjecture-for-powers-September-30-2026/build/paper.tex\#L146}{thm:main}; Statement source: \href{https://github.com/openai/math/blob/adc7f1241b42e322a6451854ab7e4b4c146bf78a/preprints/Abelian-covers-Gale-correspondences-and-the-Hodge-conjecture-for-powers-September-30-2026/build/paper.tex\#L104}{p4:abelian-covers}.
\item OpenAI, \emph{Algebraic Kuga–Satake correspondences and Hodge conjectures on a K3 quadratic locus}, 2026-09-30; citation key \texttt{OAI:Algebraic-Kuga-Satake-correspondences-and-Hodge-conjectures-on-a-K3-quadratic-locus-September-30-2026}. \href{https://github.com/openai/math/blob/adc7f1241b42e322a6451854ab7e4b4c146bf78a/preprints/Algebraic-Kuga-Satake-correspondences-and-Hodge-conjectures-on-a-K3-quadratic-locus-September-30-2026/paper.pdf}{PDF (GitHub)}; \href{https://github.com/openai/math/raw/adc7f1241b42e322a6451854ab7e4b4c146bf78a/preprints/Algebraic-Kuga-Satake-correspondences-and-Hodge-conjectures-on-a-K3-quadratic-locus-September-30-2026/paper.pdf}{PDF (direct)}; \href{https://github.com/openai/math/blob/adc7f1241b42e322a6451854ab7e4b4c146bf78a/preprints/Algebraic-Kuga-Satake-correspondences-and-Hodge-conjectures-on-a-K3-quadratic-locus-September-30-2026/README.md}{Manuscript page and citation}; \href{https://github.com/openai/math/blob/main/preprints/Algebraic-Kuga-Satake-correspondences-and-Hodge-conjectures-on-a-K3-quadratic-locus-September-30-2026/paper.pdf}{Latest version}; Statement source: \href{https://github.com/openai/math/blob/adc7f1241b42e322a6451854ab7e4b4c146bf78a/preprints/Algebraic-Kuga-Satake-correspondences-and-Hodge-conjectures-on-a-K3-quadratic-locus-September-30-2026/build/paper.tex\#L179}{main:theorem}.
\item OpenAI, \emph{Algebraicity of Kuga–Satake Correspondences for K3 Surfaces}, 2026-10-03; citation key \texttt{OAI:Algebraicity-of-Kuga-Satake-Correspondences-for-K3-Surfaces-October-3-2026}. \href{https://github.com/openai/math/blob/adc7f1241b42e322a6451854ab7e4b4c146bf78a/preprints/Algebraicity-of-Kuga-Satake-Correspondences-for-K3-Surfaces-October-3-2026/manuscript.pdf}{PDF (GitHub)}; \href{https://github.com/openai/math/raw/adc7f1241b42e322a6451854ab7e4b4c146bf78a/preprints/Algebraicity-of-Kuga-Satake-Correspondences-for-K3-Surfaces-October-3-2026/manuscript.pdf}{PDF (direct)}; \href{https://github.com/openai/math/blob/adc7f1241b42e322a6451854ab7e4b4c146bf78a/preprints/Algebraicity-of-Kuga-Satake-Correspondences-for-K3-Surfaces-October-3-2026/README.md}{Manuscript page and citation}; \href{https://github.com/openai/math/blob/main/preprints/Algebraicity-of-Kuga-Satake-Correspondences-for-K3-Surfaces-October-3-2026/manuscript.pdf}{Latest version}; Statement source: \href{https://github.com/openai/math/blob/adc7f1241b42e322a6451854ab7e4b4c146bf78a/preprints/Algebraicity-of-Kuga-Satake-Correspondences-for-K3-Surfaces-October-3-2026/build/manuscript.tex\#L109}{thm:main}.
\item OpenAI, \emph{Algebraicity of Weil classes on split abelian eightfolds}, 2026-09-18; citation key \texttt{OAI:Algebraicity-of-Weil-classes-on-split-abelian-eightfolds-September-18-2026}. \href{https://github.com/openai/math/blob/adc7f1241b42e322a6451854ab7e4b4c146bf78a/preprints/Algebraicity-of-Weil-classes-on-split-abelian-eightfolds-September-18-2026/paper.pdf}{PDF (GitHub)}; \href{https://github.com/openai/math/raw/adc7f1241b42e322a6451854ab7e4b4c146bf78a/preprints/Algebraicity-of-Weil-classes-on-split-abelian-eightfolds-September-18-2026/paper.pdf}{PDF (direct)}; \href{https://github.com/openai/math/blob/adc7f1241b42e322a6451854ab7e4b4c146bf78a/preprints/Algebraicity-of-Weil-classes-on-split-abelian-eightfolds-September-18-2026/README.md}{Manuscript page and citation}; \href{https://github.com/openai/math/blob/main/preprints/Algebraicity-of-Weil-classes-on-split-abelian-eightfolds-September-18-2026/paper.pdf}{Latest version}; Statement source: \href{https://github.com/openai/math/blob/adc7f1241b42e322a6451854ab7e4b4c146bf78a/preprints/Algebraicity-of-Weil-classes-on-split-abelian-eightfolds-September-18-2026/build/sections/01-introduction.tex\#L51}{intro:main}.
\item OpenAI, \emph{The rational Hodge conjecture for CM abelian varieties}, 2026-09-30; citation key \texttt{OAI:The-rational-Hodge-conjecture-for-CM-abelian-varieties-September-30-2026}. \href{https://github.com/openai/math/blob/adc7f1241b42e322a6451854ab7e4b4c146bf78a/preprints/The-rational-Hodge-conjecture-for-CM-abelian-varieties-September-30-2026/paper.pdf}{PDF (GitHub)}; \href{https://github.com/openai/math/raw/adc7f1241b42e322a6451854ab7e4b4c146bf78a/preprints/The-rational-Hodge-conjecture-for-CM-abelian-varieties-September-30-2026/paper.pdf}{PDF (direct)}; \href{https://github.com/openai/math/blob/adc7f1241b42e322a6451854ab7e4b4c146bf78a/preprints/The-rational-Hodge-conjecture-for-CM-abelian-varieties-September-30-2026/README.md}{Manuscript page and citation}; \href{https://github.com/openai/math/blob/main/preprints/The-rational-Hodge-conjecture-for-CM-abelian-varieties-September-30-2026/paper.pdf}{Latest version}; Statement source: \href{https://github.com/openai/math/blob/adc7f1241b42e322a6451854ab7e4b4c146bf78a/preprints/The-rational-Hodge-conjecture-for-CM-abelian-varieties-September-30-2026/build/sections/01-introduction.tex\#L17}{thm:cm}; Statement source: \href{https://github.com/openai/math/blob/adc7f1241b42e322a6451854ab7e4b4c146bf78a/preprints/The-rational-Hodge-conjecture-for-CM-abelian-varieties-September-30-2026/build/sections/07-consequences.tex\#L154}{cc:ghc}; Statement source: \href{https://github.com/openai/math/blob/adc7f1241b42e322a6451854ab7e4b4c146bf78a/preprints/The-rational-Hodge-conjecture-for-CM-abelian-varieties-September-30-2026/build/sections/07-consequences.tex\#L213}{cc:tate}; Statement source: \href{https://github.com/openai/math/blob/adc7f1241b42e322a6451854ab7e4b4c146bf78a/preprints/The-rational-Hodge-conjecture-for-CM-abelian-varieties-September-30-2026/build/sections/07-consequences.tex\#L274}{cc:standard}.
\item OpenAI, \emph{The rational Hodge conjecture for products of K3 surfaces}, 2026-10-04; citation key \texttt{OAI:The-rational-Hodge-conjecture-for-products-of-K3-surfaces-October-4-2026}. \href{https://github.com/openai/math/blob/adc7f1241b42e322a6451854ab7e4b4c146bf78a/preprints/The-rational-Hodge-conjecture-for-products-of-K3-surfaces-October-4-2026/hodge-conjecture-products-k3.pdf}{PDF (GitHub)}; \href{https://github.com/openai/math/raw/adc7f1241b42e322a6451854ab7e4b4c146bf78a/preprints/The-rational-Hodge-conjecture-for-products-of-K3-surfaces-October-4-2026/hodge-conjecture-products-k3.pdf}{PDF (direct)}; \href{https://github.com/openai/math/blob/adc7f1241b42e322a6451854ab7e4b4c146bf78a/preprints/The-rational-Hodge-conjecture-for-products-of-K3-surfaces-October-4-2026/README.md}{Manuscript page and citation}; \href{https://github.com/openai/math/blob/main/preprints/The-rational-Hodge-conjecture-for-products-of-K3-surfaces-October-4-2026/hodge-conjecture-products-k3.pdf}{Latest version}; Statement source: \href{https://github.com/openai/math/blob/adc7f1241b42e322a6451854ab7e4b4c146bf78a/preprints/The-rational-Hodge-conjecture-for-products-of-K3-surfaces-October-4-2026/build/manuscript/introduction.tex\#L14}{thm:main}.
\item OpenAI, \emph{Weil classes and Hodge classes on abelian powers}, 2026-09-30; citation key \texttt{OAI:Weil-classes-and-Hodge-classes-on-abelian-powers-September-30-2026}. \href{https://github.com/openai/math/blob/adc7f1241b42e322a6451854ab7e4b4c146bf78a/preprints/Weil-classes-and-Hodge-classes-on-abelian-powers-September-30-2026/paper.pdf}{PDF (GitHub)}; \href{https://github.com/openai/math/raw/adc7f1241b42e322a6451854ab7e4b4c146bf78a/preprints/Weil-classes-and-Hodge-classes-on-abelian-powers-September-30-2026/paper.pdf}{PDF (direct)}; \href{https://github.com/openai/math/blob/adc7f1241b42e322a6451854ab7e4b4c146bf78a/preprints/Weil-classes-and-Hodge-classes-on-abelian-powers-September-30-2026/README.md}{Manuscript page and citation}; \href{https://github.com/openai/math/blob/main/preprints/Weil-classes-and-Hodge-classes-on-abelian-powers-September-30-2026/paper.pdf}{Latest version}; Statement source: \href{https://github.com/openai/math/blob/adc7f1241b42e322a6451854ab7e4b4c146bf78a/preprints/Weil-classes-and-Hodge-classes-on-abelian-powers-September-30-2026/build/paper.tex\#L79}{p3ab:split-sixfolds}; Statement source: \href{https://github.com/openai/math/blob/adc7f1241b42e322a6451854ab7e4b4c146bf78a/preprints/Weil-classes-and-Hodge-classes-on-abelian-powers-September-30-2026/build/paper.tex\#L88}{p3ab:quadratic-five}.
\item OpenAI, \emph{Bloch's conjecture for surfaces with p\_g=0}, 2026-09-24; citation key \texttt{OAI:Blochs-Conjecture-for-Surfaces-with-pg-equals-q-equals-0-September-24-2026}. \href{https://github.com/openai/math/blob/adc7f1241b42e322a6451854ab7e4b4c146bf78a/preprints/Blochs-Conjecture-for-Surfaces-with-pg-equals-q-equals-0-September-24-2026/paper.pdf}{PDF (GitHub)}; \href{https://github.com/openai/math/raw/adc7f1241b42e322a6451854ab7e4b4c146bf78a/preprints/Blochs-Conjecture-for-Surfaces-with-pg-equals-q-equals-0-September-24-2026/paper.pdf}{PDF (direct)}; \href{https://github.com/openai/math/blob/adc7f1241b42e322a6451854ab7e4b4c146bf78a/preprints/Blochs-Conjecture-for-Surfaces-with-pg-equals-q-equals-0-September-24-2026/README.md}{Manuscript page and citation}; \href{https://github.com/openai/math/blob/main/preprints/Blochs-Conjecture-for-Surfaces-with-pg-equals-q-equals-0-September-24-2026/paper.pdf}{Latest version}; Statement source: \href{https://github.com/openai/math/blob/adc7f1241b42e322a6451854ab7e4b4c146bf78a/preprints/Blochs-Conjecture-for-Surfaces-with-pg-equals-q-equals-0-September-24-2026/build/sections/introduction.tex\#L234}{cor:tate-motives}; Statement source: \href{https://github.com/openai/math/blob/adc7f1241b42e322a6451854ab7e4b4c146bf78a/preprints/Blochs-Conjecture-for-Surfaces-with-pg-equals-q-equals-0-September-24-2026/build/sections/introduction.tex\#L18}{thm:bloch}; Statement source: \href{https://github.com/openai/math/blob/adc7f1241b42e322a6451854ab7e4b4c146bf78a/preprints/Blochs-Conjecture-for-Surfaces-with-pg-equals-q-equals-0-September-24-2026/build/sections/introduction.tex\#L31}{cor:bloch-pg-zero}.
\item \href{https://github.com/openai/math/blob/adc7f1241b42e322a6451854ab7e4b4c146bf78a/CONTENTS.md}{OpenAI catalogue entry}.
\end{itemize}
\end{document}
